\begin{lemma}
Let \(V,M,t\) be natural numbers and let \(q\ge 0\) be real.  If \(M>0\), then
\[
  \binom Vt \left(\frac{q}{M}\right)^t
  \le
  \left(\frac{eVq}{tM}\right)^t .
\]
For \(t=0\), the displayed inequality is interpreted in the Lean sense: both
sides are zeroth powers, so it asserts \(1\le 1\), even though the quotient
notation contains the factor \(tM\).
\end{lemma}
\begin{proof}
If \(t=0\), then \(\binom V0=1\), \((q/M)^0=1\), and the right-hand side is
also a zeroth power, so the desired inequality is \(1\le 1\).

Assume now that \(t>0\).  Then the denominators displayed below are nonzero.
By \noderef{OppositeProjectionChooseExponentialBound},
\[
  \binom Vt \le \left(\frac{eV}{t}\right)^t .
\]
Since \(q/M\ge 0\), its \(t\)-th power is nonnegative, so multiplying the
preceding inequality by \((q/M)^t\) preserves the inequality.  The right-hand
side is
\[
  \left(\frac{eV}{t}\right)^t \left(\frac{q}{M}\right)^t
  =
  \left(\frac{eVq}{tM}\right)^t ,
\]
which proves the claim.
\end{proof}
